\documentclass[preview]{standalone}

\usepackage{amsmath}
\usepackage{amssymb}
\usepackage{stellar}
\usepackage{definitions}
\usepackage{tikz}
\usetikzlibrary{cd}
\usepackage{bettelini}

\begin{document}

\id{sequences-functions-ascoli-arzela}
\genpage

\section{Equicontinuity}

\begin{snippetdefinition}{function-sequence-equicontinuity-definition}{Equicontinuity}
    Let \(S\subseteq\realnumbers\) and let
    \(\{f_n\}_{n\in\naturalnumbers}\) be a \sequence of \function[functions]
    \(f_n\colon S\fromto\realnumbers\). The sequence is \emph{equicontinuous}
    on \(S\) if
    \[
        \forall\varepsilon>0,\ \exists\delta>0\ \suchthat
        \forall x,y\in S,\ \forall n\in\naturalnumbers,
        \quad |x-y|<\delta\implies|f_n(x)-f_n(y)|<\varepsilon.
    \]
\end{snippetdefinition}

\begin{snippet}{function-sequence-equicontinuity-interpretation}
    Equicontinuity is uniform continuity with a modulus shared by the whole
    family. It implies that every \(f_n\) is uniformly continuous on \(S\),
    but the converse need not hold: the modulus for the individual functions
    may depend on \(n\).
\end{snippet}

\begin{snippetproposition}{uniform-lipschitz-equicontinuity-criterion}{Uniform Lipschitz criterion for equicontinuity}
    Let \(\{f_n\}_{n\in\naturalnumbers}\) be a \sequence of
    \function[functions] \(f_n\colon S\fromto\realnumbers\). If there exists
    \(L>0\), independent of \(n\), such that
    \[
        |f_n(x)-f_n(y)|\leq L|x-y|,
        \qquad \forall x,y\in S,\ \forall n\in\naturalnumbers,
    \]
    then \(\{f_n\}\) is equicontinuous on \(S\).
\end{snippetproposition}

\begin{snippetproof}{uniform-lipschitz-equicontinuity-criterion-proof}{uniform-lipschitz-equicontinuity-criterion}{Uniform Lipschitz criterion for equicontinuity}
    Given \(\varepsilon>0\), choose \(\delta=\varepsilon/L\). Then
    \(|x-y|<\delta\) implies
    \[
        |f_n(x)-f_n(y)|\leq L|x-y|<\varepsilon
    \]
    for every \(n\in\naturalnumbers\).
\end{snippetproof}

\begin{snippetcorollary}{uniformly-bounded-derivatives-equicontinuity-criterion}{Uniformly bounded derivatives imply equicontinuity}
    Let \(I\subseteq\realnumbers\) be an interval and let
    \(f_n\colon I\fromto\realnumbers\) be differentiable for every
    \(n\in\naturalnumbers\). If there exists \(L>0\) such that
    \[
        |f_n'(x)|\leq L,
        \qquad \forall x\in I,\ \forall n\in\naturalnumbers,
    \]
    then \(\{f_n\}_{n\in\naturalnumbers}\) is equicontinuous on \(I\).
\end{snippetcorollary}

\begin{snippetproof}{uniformly-bounded-derivatives-equicontinuity-criterion-proof}{uniformly-bounded-derivatives-equicontinuity-criterion}{Uniformly bounded derivatives imply equicontinuity}
    The mean value theorem gives
    \(|f_n(x)-f_n(y)|\leq L|x-y|\) for every \(x,y\in I\). The uniform
    Lipschitz criterion applies.
\end{snippetproof}

\begin{snippetproposition}{uniformly-convergent-continuous-sequence-equicontinuity-proposition}{Uniformly convergent continuous sequences are equicontinuous on compact sets}
    Let \(K\subseteq\realnumbers\) be compact and let
    \(\{f_n\}_{n\in\naturalnumbers}\) be a \sequence of continuous
    \function[functions] \(f_n\colon K\fromto\realnumbers\). If \(f_n\)
    converges uniformly on \(K\), then \(\{f_n\}\) is equicontinuous on
    \(K\).
\end{snippetproposition}

\begin{snippetproof}{uniformly-convergent-continuous-sequence-equicontinuity-proposition-proof}{uniformly-convergent-continuous-sequence-equicontinuity-proposition}{Uniformly convergent continuous sequences are equicontinuous on compact sets}
    Let \(f_n\to f\) uniformly and fix \(\varepsilon>0\). The function \(f\)
    is continuous, hence uniformly continuous on the compact set \(K\). Choose
    \(N\) such that
    \[
        |f_n(x)-f(x)|<\frac{\varepsilon}{3},
        \qquad \forall n>N,\ \forall x\in K,
    \]
    and choose \(\delta_0>0\) such that
    \[
        |x-y|<\delta_0\implies
        |f(x)-f(y)|<\frac{\varepsilon}{3}.
    \]
    Then, for \(n>N\) and \(|x-y|<\delta_0\),
    \begin{align*}
        |f_n(x)-f_n(y)|
        &\leq |f_n(x)-f(x)|+|f(x)-f(y)|+|f(y)-f_n(y)|\\
        &<\varepsilon.
    \end{align*}
    Each of the finitely many functions \(f_1,\ldots,f_N\) is uniformly
    continuous on \(K\). Let \(\delta_i>0\) be a suitable modulus for
    \(f_i\), and take
    \[
        \delta=\min\{\delta_0,\delta_1,\ldots,\delta_N\}.
    \]
    This \(\delta\) works for every \(n\in\naturalnumbers\).
\end{snippetproof}

\section{From pointwise to uniform convergence}

\begin{snippetlemma}{equicontinuous-pointwise-convergence-uniform-lemma}{Equicontinuous pointwise convergence on a compact set}
    Let \(K\subseteq\realnumbers\) be compact and let
    \(\{f_n\}_{n\in\naturalnumbers}\) be an equicontinuous \sequence of
    \function[functions] \(f_n\colon K\fromto\realnumbers\). If \(f_n\)
    converges pointwise on \(K\) to \(f\colon K\fromto\realnumbers\), then
    \(f_n\) converges uniformly to \(f\) on \(K\). In particular, \(f\) is
    continuous.
\end{snippetlemma}

\begin{snippetproof}{equicontinuous-pointwise-convergence-uniform-lemma-proof}{equicontinuous-pointwise-convergence-uniform-lemma}{Equicontinuous pointwise convergence on a compact set}
    Fix \(\varepsilon>0\). Equicontinuity yields \(\delta>0\) such that
    \[
        |x-y|<\delta\implies
        |f_i(x)-f_i(y)|<\frac{\varepsilon}{3}
    \]
    for every \(i\in\naturalnumbers\). Compactness provides points
    \(x_1,\ldots,x_r\in K\) such that
    \[
        K\subseteq\bigcup_{j=1}^{r}B_\delta(x_j).
    \]
    At each \(x_j\), the numerical sequence \(\{f_n(x_j)\}\) is Cauchy.
    Hence there exists \(N_j\) such that
    \[
        |f_n(x_j)-f_m(x_j)|<\frac{\varepsilon}{3}
    \]
    whenever \(m,n>N_j\). Let \(N=\max\{N_1,\ldots,N_r\}\). Given
    \(x\in K\), choose \(x_j\) with \(|x-x_j|<\delta\). For every
    \(m,n>N\),
    \begin{align*}
        |f_n(x)-f_m(x)|
        &\leq |f_n(x)-f_n(x_j)|+|f_n(x_j)-f_m(x_j)|
        +|f_m(x_j)-f_m(x)|\\
        &<\varepsilon.
    \end{align*}
    Thus \(\{f_n\}\) is uniformly Cauchy and therefore converges uniformly.
    Its uniform limit must equal the pointwise limit \(f\).
\end{snippetproof}

\section{Ascoli--Arzelà theorem}

\begin{snippettheorem}{ascoli-arzela-sequence-theorem}{Ascoli--Arzelà theorem}
    Let \(K\subseteq\realnumbers\) be compact and let
    \(\{f_n\}_{n\in\naturalnumbers}\) be a \sequence of \function[functions]
    \(f_n\colon K\fromto\realnumbers\). If the sequence is equicontinuous and
    uniformly bounded on \(K\), then it has a subsequence that converges
    uniformly on \(K\).
\end{snippettheorem}

\begin{snippetproof}{ascoli-arzela-sequence-theorem-proof}{ascoli-arzela-sequence-theorem}{Ascoli--Arzelà theorem}
    Choose a countable dense subset
    \(\{r_j\}_{j\in\naturalnumbers}\subseteq K\). Uniform boundedness makes
    the numerical sequence \(\{f_n(r_1)\}\) bounded, so Bolzano--Weierstrass
    gives a convergent subsequence, denoted \(\{f_n^1\}\). From it extract a
    subsequence \(\{f_n^2\}\) that converges at \(r_2\). Continuing gives
    nested subsequences
    \[
        \{f_n\}\supseteq\{f_n^1\}\supseteq\{f_n^2\}
        \supseteq\{f_n^3\}\supseteq\cdots,
    \]
    where \(\{f_n^j(r_j)\}\) converges for every \(j\).

    Consider the diagonal subsequence \(g_k=f_k^k\):
    \begin{center}
        \begin{tikzcd}
            f_1^1 \arrow[rd] & f_2^1            & f_3^1            & f_4^1  & \cdots \\
            f_1^2            & f_2^2 \arrow[rd] & f_3^2            & f_4^2  & \cdots \\
            f_1^3            & f_2^3            & f_3^3 \arrow[rd] & f_4^3  & \cdots \\
            f_1^4            & f_2^4            & f_3^4            & f_4^4  & \cdots \\
            \vdots           & \vdots           & \vdots           & \vdots & \ddots
        \end{tikzcd}
    \end{center}
    By construction, \(\{g_k(r_j)\}\) converges for every \(j\).

    We show that \(\{g_k(x)\}\) is Cauchy for every \(x\in K\). Fix
    \(\varepsilon>0\). Equicontinuity gives \(\delta>0\) such that
    \[
        |g_k(x)-g_k(y)|<\frac{\varepsilon}{3}
    \]
    for every \(k\) whenever \(|x-y|<\delta\). By density, choose
    \(r_j\) with \(|x-r_j|<\delta\). Since \(\{g_k(r_j)\}\) is Cauchy,
    there exists \(N=N(\varepsilon,x)\) such that
    \[
        |g_n(r_j)-g_m(r_j)|<\frac{\varepsilon}{3}
    \]
    for \(m,n>N\). Therefore,
    \begin{align*}
        |g_n(x)-g_m(x)|
        &\leq |g_n(x)-g_n(r_j)|+|g_n(r_j)-g_m(r_j)|
        +|g_m(r_j)-g_m(x)|\\
        &<\varepsilon.
    \end{align*}
    Thus \(\{g_k\}\) converges pointwise on \(K\). It remains
    equicontinuous as a subsequence of \(\{f_n\}\), so the preceding lemma
    upgrades the convergence to uniform convergence.
\end{snippetproof}

\begin{snippetexample}{uniformly-bounded-sequence-without-uniformly-convergent-subsequence-example}{Uniform boundedness alone is insufficient}
    On \(K=[0,1]\), let
    \[
        f_n(x)=\frac{x^2}{x^2+(1-nx)^2}.
    \]
    The sequence converges pointwise to zero and satisfies
    \(|f_n(x)|\leq1\), but \(f_n(1/n)=1\). Every pointwise convergent
    subsequence must therefore have pointwise limit zero, whereas for any
    subsequence \(\{f_{n_k}\}\),
    \[
        \sup_{x\in[0,1]}|f_{n_k}(x)|
        \geq f_{n_k}\left(\frac1{n_k}\right)=1.
    \]
    Hence no subsequence converges uniformly. In particular, the original
    family cannot be equicontinuous.
\end{snippetexample}

\end{document}
